\documentclass[11pt]{article}

\usepackage[utf8]{inputenc}
\usepackage[T1]{fontenc}
\usepackage[spanish,es-nodecimaldot]{babel}
\usepackage{amsmath,amssymb,mathtools}
\usepackage{geometry}
\usepackage{booktabs}
\usepackage{enumitem}
\usepackage{array}
\usepackage{hyperref}
\usepackage{xcolor}
\usepackage{siunitx}

\geometry{margin=2.5cm}
\hypersetup{
    colorlinks=true,
    linkcolor=blue!50!black,
    urlcolor=blue!50!black
}

\title{Modelado Mec\'anico y Linealizaci\'on Rigurosa\\Plataforma Antivibratoria Activa}
\author{Versi\'on detallada para sustentar en clase y pasar a Overleaf}
\date{15 de abril de 2026}

\begin{document}

\maketitle
\tableofcontents
\newpage

\section{Objetivo de este documento}

La idea aqu\'i no es solo escribir ecuaciones, sino dejar completamente claro:

\begin{itemize}[leftmargin=1.2cm]
\item qu\'e ecuaciones salen de \textbf{leyes f\'isicas fundamentales};
\item qu\'e ecuaciones salen de \textbf{relaciones constitutivas} de ingenier\'ia;
\item qu\'e t\'erminos son \textbf{no linealidades realistas};
\item c\'omo se construye el \textbf{modelo no lineal};
\item c\'omo se obtiene el \textbf{modelo linealizado};
\item qu\'e expresiones se simplifican, por qu\'e se simplifican y en qu\'e paso.
\end{itemize}

La exigencia de ``m\'as rigor'' no significa eliminar todos los supuestos, porque eso es imposible en un modelo ingenieril. Significa otra cosa:

\begin{quote}
\textbf{cada supuesto debe estar escrito, debe ser localizable en una etapa concreta del modelado y debe tener justificaci\'on f\'isica o metodol\'ogica.}
\end{quote}

\section{Clasificaci\'on correcta de las ecuaciones que vamos a usar}

Una parte importante del rigor es no mezclar leyes fundamentales con aproximaciones emp\'iricas. En este proyecto aparecen tres niveles:

\subsection{Leyes f\'isicas fundamentales}

\begin{itemize}[leftmargin=1.2cm]
\item \textbf{Segunda ley de Newton:}
\[
\sum F = m a.
\]
\item \textbf{Ley de voltajes de Kirchhoff (KVL):}
\[
\sum v = 0.
\]
\item \textbf{Ley de Lorentz / actuador electromagn\'etico lineal:} en un actuador tipo voice coil la fuerza es proporcional a la corriente.
\item \textbf{Ley de Faraday:} el movimiento del actuador produce una fuerza contraelectromotriz proporcional a la velocidad relativa.
\end{itemize}

\subsection{Relaciones constitutivas de ingenier\'ia}

No son leyes universales de la naturaleza, sino modelos aceptados en un rango de operaci\'on:

\begin{itemize}[leftmargin=1.2cm]
\item \textbf{Ley de Hooke:}
\[
F_s = -k\,\Delta \ell,
\]
v\'alida para deformaciones peque\~nas alrededor de la configuraci\'on de equilibrio.
\item \textbf{Amortiguamiento viscoso ideal:}
\[
F_d = -c\,\Delta \dot{\ell}.
\]
\end{itemize}

\subsection{Regularizaciones y efectos de realismo}

Estos t\'erminos no son leyes fundamentales. Se a\~naden para representar mejor una planta real:

\begin{itemize}[leftmargin=1.2cm]
\item rigidez c\'ubica;
\item fricci\'on seca regularizada con $\tanh$;
\item saturaci\'on del actuador;
\item topes mec\'anicos.
\end{itemize}

Esto es importante para sustentar bien el proyecto: si te preguntan de d\'onde sale un t\'ermino, debes poder decir \textbf{si es ley f\'isica, relaci\'on constitutiva o regularizaci\'on de simulaci\'on}.

\section{Sistema f\'isico que se modela}

Se dise\~na una plataforma antivibratoria activa formada por:

\begin{itemize}[leftmargin=1.2cm]
\item una base inferior que se mueve con una perturbaci\'on externa;
\item una plataforma superior donde ir\'ia montado un equipo sensible;
\item una uni\'on el\'astica entre base y plataforma;
\item una disipaci\'on mec\'anica entre base y plataforma;
\item un actuador lineal que aplica fuerza entre base y plataforma.
\end{itemize}

La variable que se quiere reducir es el movimiento de la plataforma superior.

\section{Variables, coordenadas y convenci\'on de signos}

\subsection{Coordenadas absolutas}

Definimos:
\begin{equation}
x(t) = \text{posici\'on absoluta de la plataforma},
\end{equation}
\begin{equation}
z(t) = \text{posici\'on absoluta de la base}.
\end{equation}

\subsection{Coordenada relativa}

Definimos la deformaci\'on relativa entre plataforma y base:
\begin{equation}
q(t) = x(t) - z(t).
\label{eq:q}
\end{equation}

Entonces:
\begin{equation}
\dot q(t) = \dot x(t) - \dot z(t),
\end{equation}
\begin{equation}
\ddot q(t) = \ddot x(t) - \ddot z(t).
\end{equation}

\subsection{Por qu\'e esta coordenada es la correcta}

El resorte y el amortiguador no ``ven'' la posici\'on absoluta de la plataforma, sino \textbf{cu\'anto se estiran o se comprimen}. Esa deformaci\'on es exactamente $q=x-z$.

Por eso, si se escribe un t\'ermino como $kx$ sin m\'as contexto, en este sistema estar\'ia mal planteado. El t\'ermino correcto es $k(x-z)=kq$.

\section{Sobre la gravedad: c\'omo tratarla sin esconderla}

Aqu\'i suele aparecer una duda correcta: \textit{si la plataforma tiene masa, por qu\'e no aparece el peso $mg$?}

La respuesta depende del eje elegido.

\subsection{Caso 1: eje de movimiento horizontal}

Si la vibraci\'on analizada es horizontal, la gravedad no tiene componente sobre el eje de modelado. En ese caso no aparece en la ecuaci\'on din\'amica sobre ese eje.

\subsection{Caso 2: eje de movimiento vertical}

Si el movimiento fuese vertical, entonces s\'i aparece el peso. En coordenadas absolutas:
\begin{equation}
m\ddot x = -k(x-z) - c(\dot x-\dot z) + F_a - mg.
\label{eq:vertical_raw}
\end{equation}

Sin embargo, en control normalmente se trabaja alrededor del equilibrio est\'atico. Si el equilibrio est\'atico satisface:
\begin{equation}
0 = -k(x_0-z_0) + F_{a0} - mg,
\label{eq:staticeq}
\end{equation}
y se definen variables incrementales:
\[
\tilde x = x-x_0,\qquad
\tilde z = z-z_0,\qquad
\tilde F_a = F_a-F_{a0},
\]
entonces al reemplazar en \eqref{eq:vertical_raw} y usar \eqref{eq:staticeq}, los t\'erminos constantes se cancelan y queda:
\begin{equation}
m\ddot{\tilde x} = -k(\tilde x-\tilde z)-c(\dot{\tilde x}-\dot{\tilde z})+\tilde F_a.
\end{equation}

Es decir, \textbf{la gravedad no desaparece por magia}; desaparece porque se absorbe en el equilibrio est\'atico y luego se modelan desviaciones respecto a ese equilibrio.

En lo que sigue, para no recargar la notaci\'on, trabajaremos directamente con variables incrementales y volveremos a escribirlas como $x$, $z$, $q$ e $i$.

\section{Derivaci\'on mec\'anica desde la segunda ley de Newton}

\subsection{Fuerzas que act\'uan sobre la plataforma}

Sobre la masa de la plataforma act\'uan:

\begin{itemize}[leftmargin=1.2cm]
\item fuerza el\'astica del resorte;
\item fuerza viscosa del amortiguador;
\item fuerza activa del actuador.
\end{itemize}

Aplicando la segunda ley de Newton:
\begin{equation}
\sum F = m\ddot x.
\label{eq:newton}
\end{equation}

\subsection{Fuerza del resorte}

Por ley de Hooke, para deformaciones peque\~nas alrededor del punto de operaci\'on:
\begin{equation}
F_s = -k(x-z) = -kq.
\label{eq:spring}
\end{equation}

El signo negativo indica que el resorte se opone a la deformaci\'on.

\subsection{Fuerza del amortiguador}

Con el modelo viscoso ideal:
\begin{equation}
F_d = -c(\dot x-\dot z) = -c\dot q.
\label{eq:damper}
\end{equation}

El signo negativo indica que el amortiguador se opone a la velocidad relativa.

\subsection{Fuerza del actuador}

La fuerza del actuador se denota por $F_a$. Tomamos como convenci\'on que $F_a>0$ empuja la plataforma en el sentido positivo del eje.

\subsection{Ecuaci\'on mec\'anica en coordenadas absolutas}

Sumando fuerzas:
\begin{equation}
m\ddot x = -k(x-z)-c(\dot x-\dot z)+F_a.
\label{eq:mech_abs}
\end{equation}

\subsection{Paso a coordenada relativa}

Usando $x=q+z$, $\dot x=\dot q+\dot z$ y $\ddot x=\ddot q+\ddot z$, la ecuaci\'on \eqref{eq:mech_abs} se transforma en:
\begin{equation}
m(\ddot q+\ddot z) = -kq-c\dot q+F_a.
\end{equation}

Reordenando:
\begin{equation}
m\ddot q + c\dot q + kq = F_a - m\ddot z.
\label{eq:mech_rel}
\end{equation}

\subsection{Interpretaci\'on de la perturbaci\'on}

La base no entra como fuerza directa, sino como aceleraci\'on de soporte. En coordenada relativa eso aparece como una fuerza inercial equivalente:
\begin{equation}
F_{\text{pert}} = -m\ddot z.
\end{equation}

Ese punto es importante porque te permite defender por qu\'e la perturbaci\'on de base entra con signo negativo y multiplicada por la masa.

\section{Derivaci\'on rigurosa del actuador electromec\'anico}

Aqu\'i es donde se obtiene el \textbf{tercer estado} del sistema sin inventar una din\'amica artificial.

\subsection{Fuerza electromagn\'etica}

En un actuador lineal tipo voice coil o equivalente, para un rango de operaci\'on alrededor del punto de trabajo, la fuerza es proporcional a la corriente:
\begin{equation}
F_a = K_t i.
\label{eq:Kt}
\end{equation}

Esta relaci\'on proviene de la interacci\'on entre corriente y campo magn\'etico. En un modelo de ingenier\'ia de bajo orden, toda esa f\'isica se condensa en la constante $K_t$.

\subsection{Ecuaci\'on el\'ectrica por KVL}

El circuito del actuador se modela con una resistencia $R$, una inductancia $L$ y una fuerza contraelectromotriz $e_b$. Por KVL:
\begin{equation}
u = L\dot i + Ri + e_b.
\label{eq:kvl}
\end{equation}

\subsection{Back-EMF por movimiento relativo}

La ley de Faraday indica que el movimiento de la parte m\'ovil del actuador induce un voltaje. Como el actuador est\'a entre la base y la plataforma, la velocidad relevante es la velocidad relativa:
\begin{equation}
e_b = K_e \dot q.
\label{eq:bemf}
\end{equation}

\subsection{Modelo el\'ectrico del actuador}

Sustituyendo \eqref{eq:bemf} en \eqref{eq:kvl}:
\begin{equation}
L\dot i + Ri + K_e\dot q = u.
\label{eq:elec}
\end{equation}

\subsection{Comentario de rigor}

Esta ecuaci\'on no es una suposici\'on arbitraria. Lo que se est\'a suponiendo, de manera expl\'icita, es:

\begin{itemize}[leftmargin=1.2cm]
\item campo magn\'etico aproximadamente constante en la zona de operaci\'on;
\item linealidad local entre fuerza y corriente;
\item linealidad local entre back-EMF y velocidad relativa;
\item par\'ametros $R$, $L$, $K_t$ y $K_e$ aproximadamente constantes.
\end{itemize}

Esas son suposiciones de ingenier\'ia razonables y defendibles en un proyecto de control de bajo orden.

\section{Modelo electrome\'canico acoplado exacto de bajo orden}

Sustituyendo \eqref{eq:Kt} en \eqref{eq:mech_rel} se obtiene:
\begin{equation}
m\ddot q + c\dot q + kq = K_t i - m\ddot z.
\label{eq:mech_exact}
\end{equation}

Junto con \eqref{eq:elec}:
\begin{equation}
L\dot i + Ri + K_e\dot q = u.
\label{eq:elec_exact}
\end{equation}

Este es el \textbf{modelo lineal exacto de tercer orden} del sistema de bajo orden. El orden total es tres porque aparecen:

\begin{itemize}[leftmargin=1.2cm]
\item una coordenada mec\'anica $q$;
\item su velocidad $\dot q$;
\item una variable el\'ectrica interna $i$.
\end{itemize}

\section{Forma en espacio de estados}

Definimos:
\[
x_1=q,\qquad x_2=\dot q,\qquad x_3=i.
\]

Tambi\'en definimos la perturbaci\'on:
\[
w(t)=\ddot z(t).
\]

Entonces:
\begin{equation}
\dot x_1 = x_2,
\end{equation}
\begin{equation}
\dot x_2 = -\frac{k}{m}x_1-\frac{c}{m}x_2+\frac{K_t}{m}x_3-w,
\end{equation}
\begin{equation}
\dot x_3 = -\frac{K_e}{L}x_2-\frac{R}{L}x_3+\frac{1}{L}u.
\end{equation}

En forma matricial:
\begin{equation}
\dot x =
\underbrace{
\begin{bmatrix}
0 & 1 & 0\\
-\dfrac{k}{m} & -\dfrac{c}{m} & \dfrac{K_t}{m}\\
0 & -\dfrac{K_e}{L} & -\dfrac{R}{L}
\end{bmatrix}}_{A}
x
+
\underbrace{
\begin{bmatrix}
0\\
0\\
\dfrac{1}{L}
\end{bmatrix}}_{B}
u
+
\underbrace{
\begin{bmatrix}
0\\
-1\\
0
\end{bmatrix}}_{E}
w.
\label{eq:ss_exact}
\end{equation}

Si la salida es la deformaci\'on relativa:
\begin{equation}
y=q=
\underbrace{
\begin{bmatrix}
1 & 0 & 0
\end{bmatrix}}_{C}
x.
\end{equation}

Si la salida es la posici\'on absoluta de la plataforma, entonces:
\begin{equation}
y=x=q+z.
\end{equation}

\section{Modelo no lineal realista}

Para construir el modelo realista de simulaci\'on se agregan efectos que no queremos necesariamente en el modelo nominal de dise\~no.

\subsection{Rigidez no lineal}

Si la ley fuerza-deformaci\'on del elemento el\'astico deja de ser perfectamente lineal, una aproximaci\'on sim\'etrica de primer orden no lineal es:
\begin{equation}
F_{s,\text{nl}} = -kq - k_3 q^3.
\label{eq:cubic}
\end{equation}

El t\'ermino c\'ubico aparece porque el primer t\'ermino no lineal impar compatible con una simetr\'ia de restauraci\'on es $q^3$.

\subsection{Fricci\'on seca regularizada}

La fricci\'on de Coulomb ideal es no diferenciable en velocidad cero. Para poder simular y linealizar localmente, se usa una regularizaci\'on:
\begin{equation}
F_f(\dot q) = F_c \tanh\left(\frac{\dot q}{v_\varepsilon}\right).
\label{eq:fric}
\end{equation}

Esto no es una ley fundamental. Es una aproximaci\'on suave que:

\begin{itemize}[leftmargin=1.2cm]
\item se comporta como fricci\'on seca para $|\dot q| \gg v_\varepsilon$;
\item es diferenciable en $\dot q=0$;
\item permite obtener un Jacobiano.
\end{itemize}

\subsection{Saturaci\'on del actuador}

En la electr\'onica real no se puede aplicar cualquier voltaje:
\begin{equation}
u_{\text{act}} = \operatorname{sat}(u).
\label{eq:sat}
\end{equation}

\subsection{Topes mec\'anicos}

Si la deformaci\'on supera un recorrido permitido $q_{\max}$, aparece una fuerza de contacto. Un modelo de ingenier\'ia sencillo es:
\begin{equation}
F_{\text{topes}}(q,\dot q)=
\begin{cases}
-k_t(q-q_{\max})-c_t\dot q, & q>q_{\max},\\[4pt]
0, & |q|\le q_{\max},\\[4pt]
-k_t(q+q_{\max})-c_t\dot q, & q<-q_{\max}.
\end{cases}
\label{eq:stops}
\end{equation}

\subsection{Modelo completo no lineal}

Con todo lo anterior:
\begin{equation}
m\ddot q =
-kq
-c\dot q
-k_3 q^3
-F_c \tanh\left(\frac{\dot q}{v_\varepsilon}\right)
+F_{\text{topes}}(q,\dot q)
+K_t i
-m\ddot z,
\label{eq:nonlin_mech}
\end{equation}
\begin{equation}
L\dot i + Ri + K_e\dot q = \operatorname{sat}(u).
\label{eq:nonlin_elec}
\end{equation}

Estas dos ecuaciones representan el \textbf{modelo realista} que luego se quiere simular en Simulink/Simscape.

\section{Estado no lineal y punto de equilibrio}

Definimos:
\[
x=
\begin{bmatrix}
q\\
v\\
i
\end{bmatrix},
\qquad v=\dot q,
\qquad w=\ddot z.
\]

Entonces:
\begin{equation}
\dot x_1 = x_2,
\end{equation}
\begin{equation}
\dot x_2 =
\frac{1}{m}
\left(
-kx_1
-cx_2
-k_3 x_1^3
-F_c\tanh\left(\frac{x_2}{v_\varepsilon}\right)
+F_{\text{topes}}(x_1,x_2)
+K_t x_3
\right)
-w,
\label{eq:f2}
\end{equation}
\begin{equation}
\dot x_3 =
\frac{1}{L}
\left(
-K_e x_2
-R x_3
+\operatorname{sat}(u)
\right).
\label{eq:f3}
\end{equation}

El punto natural de linealizaci\'on es:
\begin{equation}
x^\star=
\begin{bmatrix}
0\\
0\\
0
\end{bmatrix},
\qquad
u^\star=0,
\qquad
w^\star=0.
\label{eq:eqpoint}
\end{equation}

F\'isicamente, este equilibrio representa:

\begin{itemize}[leftmargin=1.2cm]
\item plataforma centrada respecto a la base;
\item velocidad relativa nula;
\item corriente nula;
\item base sin aceleraci\'on externa;
\item actuador trabajando dentro de la zona no saturada;
\item topes inactivos.
\end{itemize}

\section{Linealizaci\'on rigurosa por Taylor de primer orden}

\subsection{Idea general}

Si el sistema no lineal se escribe como:
\[
\dot x = f(x,u,w),
\]
entonces alrededor del equilibrio \eqref{eq:eqpoint} se definen desviaciones:
\[
\delta x = x-x^\star,\qquad
\delta u = u-u^\star,\qquad
\delta w = w-w^\star.
\]

La expansi\'on de Taylor de primer orden da:
\begin{equation}
\delta \dot x
\approx
\left.\frac{\partial f}{\partial x}\right|_\star \delta x
+
\left.\frac{\partial f}{\partial u}\right|_\star \delta u
+
\left.\frac{\partial f}{\partial w}\right|_\star \delta w.
\label{eq:taylor}
\end{equation}

Definimos:
\begin{equation}
A=\left.\frac{\partial f}{\partial x}\right|_\star,\qquad
B=\left.\frac{\partial f}{\partial u}\right|_\star,\qquad
E=\left.\frac{\partial f}{\partial w}\right|_\star.
\end{equation}

\subsection{Primera ecuaci\'on}

Como
\[
f_1(x,u,w)=x_2,
\]
entonces:
\[
\frac{\partial f_1}{\partial x_1}=0,\qquad
\frac{\partial f_1}{\partial x_2}=1,\qquad
\frac{\partial f_1}{\partial x_3}=0,\qquad
\frac{\partial f_1}{\partial u}=0,\qquad
\frac{\partial f_1}{\partial w}=0.
\]

\subsection{Segunda ecuaci\'on: derivadas t\'ermino a t\'ermino}

Tomamos:
\[
f_2(x,u,w)=
\frac{1}{m}
\left(
-kx_1
-cx_2
-k_3 x_1^3
-F_c\tanh\left(\frac{x_2}{v_\varepsilon}\right)
+F_{\text{topes}}(x_1,x_2)
+K_t x_3
\right)
-w.
\]

\subsubsection{Derivada respecto a $x_1$}

Se deriva cada t\'ermino dependiente de $x_1$:
\[
\frac{\partial (-kx_1)}{\partial x_1} = -k,
\qquad
\frac{\partial (-k_3x_1^3)}{\partial x_1} = -3k_3x_1^2.
\]

Evaluando en el equilibrio $x_1^\star=0$:
\[
-3k_3(x_1^\star)^2 = 0.
\]

Adem\'as, como el equilibrio est\'a lejos del contacto con topes:
\[
\left.\frac{\partial F_{\text{topes}}}{\partial x_1}\right|_\star = 0.
\]

Por tanto:
\begin{equation}
\left.\frac{\partial f_2}{\partial x_1}\right|_\star
=
-\frac{k}{m}.
\label{eq:df2dx1}
\end{equation}

\subsubsection{Derivada respecto a $x_2$}

Primero:
\[
\frac{\partial (-cx_2)}{\partial x_2}=-c.
\]

Luego, para la fricci\'on regularizada:
\[
\frac{\partial}{\partial x_2}
\left[
-F_c\tanh\left(\frac{x_2}{v_\varepsilon}\right)
\right]
=
-F_c\frac{1}{v_\varepsilon}\operatorname{sech}^2\left(\frac{x_2}{v_\varepsilon}\right).
\]

Evaluando en $x_2^\star=0$:
\[
\operatorname{sech}^2(0)=1,
\]
por tanto:
\[
\left.
\frac{\partial}{\partial x_2}
\left[
-F_c\tanh\left(\frac{x_2}{v_\varepsilon}\right)
\right]
\right|_\star
=
-\frac{F_c}{v_\varepsilon}.
\]

Adem\'as, si no hay contacto con topes en el equilibrio:
\[
\left.\frac{\partial F_{\text{topes}}}{\partial x_2}\right|_\star = 0.
\]

Entonces:
\begin{equation}
\left.\frac{\partial f_2}{\partial x_2}\right|_\star
=
-\frac{1}{m}\left(c+\frac{F_c}{v_\varepsilon}\right).
\label{eq:df2dx2}
\end{equation}

\subsubsection{Derivada respecto a $x_3$}

Como el acoplamiento con corriente es lineal:
\begin{equation}
\left.\frac{\partial f_2}{\partial x_3}\right|_\star
=
\frac{K_t}{m}.
\label{eq:df2dx3}
\end{equation}

\subsubsection{Derivada respecto a $u$}

La segunda ecuaci\'on no depende directamente de $u$:
\begin{equation}
\left.\frac{\partial f_2}{\partial u}\right|_\star = 0.
\end{equation}

\subsubsection{Derivada respecto a $w$}

Como $f_2$ contiene $-w$:
\begin{equation}
\left.\frac{\partial f_2}{\partial w}\right|_\star = -1.
\label{eq:df2dw}
\end{equation}

\subsection{Tercera ecuaci\'on: derivadas t\'ermino a t\'ermino}

Tomamos:
\[
f_3(x,u,w)=
\frac{1}{L}
\left(
-K_e x_2
-R x_3
+\operatorname{sat}(u)
\right).
\]

\subsubsection{Derivadas respecto a los estados}

\[
\left.\frac{\partial f_3}{\partial x_1}\right|_\star=0,
\qquad
\left.\frac{\partial f_3}{\partial x_2}\right|_\star=-\frac{K_e}{L},
\qquad
\left.\frac{\partial f_3}{\partial x_3}\right|_\star=-\frac{R}{L}.
\]

\subsubsection{Derivada respecto a la entrada}

Aqu\'i hay una condici\'on importante. La saturaci\'on solo es derivable con pendiente unitaria dentro de la zona lineal. Como el equilibrio se elige en $u^\star=0$ y asumimos que $0$ est\'a dentro de la zona no saturada:
\[
\left.\frac{d\,\operatorname{sat}(u)}{du}\right|_\star = 1.
\]

Entonces:
\begin{equation}
\left.\frac{\partial f_3}{\partial u}\right|_\star = \frac{1}{L}.
\label{eq:df3du}
\end{equation}

\subsubsection{Derivada respecto a la perturbaci\'on}

La tercera ecuaci\'on no depende de $w$:
\[
\left.\frac{\partial f_3}{\partial w}\right|_\star = 0.
\]

\section{Matrices linealizadas}

Con \eqref{eq:df2dx1}--\eqref{eq:df3du}, el modelo linealizado local es:
\begin{equation}
\delta \dot x = A\,\delta x + B\,\delta u + E\,\delta w,
\end{equation}
con
\begin{equation}
A=
\begin{bmatrix}
0 & 1 & 0\\[4pt]
-\dfrac{k}{m} & -\dfrac{1}{m}\left(c+\dfrac{F_c}{v_\varepsilon}\right) & \dfrac{K_t}{m}\\[8pt]
0 & -\dfrac{K_e}{L} & -\dfrac{R}{L}
\end{bmatrix},
\qquad
B=
\begin{bmatrix}
0\\
0\\
\dfrac{1}{L}
\end{bmatrix},
\qquad
E=
\begin{bmatrix}
0\\
-1\\
0
\end{bmatrix}.
\label{eq:lin_full}
\end{equation}

\section{Qu\'e expresiones se simplifican y por qu\'e}

Aqu\'i est\'a una de las partes m\'as importantes para sustentar.

\subsection{Por qu\'e el t\'ermino c\'ubico desaparece}

En el equilibrio:
\[
\frac{d}{dq}(k_3 q^3)=3k_3q^2.
\]

Como $q^\star=0$:
\[
3k_3(q^\star)^2 = 0.
\]

Eso significa que \textbf{el t\'ermino c\'ubico no contribuye en primer orden}. No se elimina por capricho, sino porque su derivada en el punto de equilibrio es cero.

\subsection{Por qu\'e los topes desaparecen}

Si el equilibrio est\'a dentro de la zona libre, entonces alrededor de ese punto:
\[
F_{\text{topes}}(q,\dot q)=0
\]
en un vecindario del equilibrio. Por eso su Jacobiano local es cero.

\subsection{Por qu\'e la saturaci\'on se vuelve lineal}

Si el punto de operaci\'on est\'a dentro de la regi\'on no saturada, entonces localmente:
\[
\operatorname{sat}(u)\approx u.
\]

Esa es una aproximaci\'on local, no global. Si el actuador opera saturado, esta linealizaci\'on deja de ser v\'alida.

\subsection{Por qu\'e la fricci\'on regularizada s\'i deja una huella lineal}

Como
\[
\tanh(\xi)\approx \xi \quad \text{para } |\xi|\ll 1,
\]
entonces, cerca de velocidad cero:
\begin{equation}
F_c\tanh\left(\frac{\dot q}{v_\varepsilon}\right)
\approx
\frac{F_c}{v_\varepsilon}\dot q.
\label{eq:fric_lin}
\end{equation}

Por eso, la linealizaci\'on estricta del modelo realista introduce un amortiguamiento equivalente adicional:
\begin{equation}
c_{\text{eq}} = c + \frac{F_c}{v_\varepsilon}.
\label{eq:ceq}
\end{equation}

\section{Dos modelos lineales distintos y ambos son defendibles}

\subsection{Modelo lineal 1: linealizaci\'on estricta del modelo realista}

Si linealizas literalmente el modelo no lineal completo, obtienes la matriz \eqref{eq:lin_full}, con amortiguamiento equivalente $c_{\text{eq}}$.

\subsection{Modelo lineal 2: modelo nominal limpio para dise\~no}

En muchos cursos de control se prefiere construir el modelo nominal para dise\~no ignorando:

\begin{itemize}[leftmargin=1.2cm]
\item fricci\'on seca;
\item topes;
\item saturaci\'on;
\item no linealidad c\'ubica.
\end{itemize}

La justificaci\'on es que esos efectos se reservan para el \textbf{modelo realista}, mientras que el dise\~no del controlador se hace sobre la din\'amica principal.

En ese caso el modelo nominal queda:
\begin{equation}
m\ddot q + c\dot q + kq = K_t i - m\ddot z,
\label{eq:nom_mech}
\end{equation}
\begin{equation}
L\dot i + Ri + K_e\dot q = u.
\label{eq:nom_elec}
\end{equation}

Este sigue siendo un modelo \textbf{de tercer orden}, pero m\'as limpio para an\'alisis.

\section{Funci\'on de transferencia del modelo nominal exacto}

Para obtener la planta de control, primero se toma perturbaci\'on nula:
\[
w=\ddot z=0.
\]

Aplicando Laplace a \eqref{eq:nom_mech} y \eqref{eq:nom_elec}, con condiciones iniciales nulas:
\[
(ms^2+cs+k)Q(s)=K_t I(s),
\]
\[
(Ls+R)I(s)+K_e sQ(s)=U(s).
\]

Despejando $I(s)$:
\[
I(s)=\frac{ms^2+cs+k}{K_t}Q(s).
\]

Sustituyendo en la ecuaci\'on el\'ectrica:
\[
(Ls+R)\frac{ms^2+cs+k}{K_t}Q(s)+K_e sQ(s)=U(s).
\]

Multiplicando por $K_t$:
\[
\left[(Ls+R)(ms^2+cs+k)+K_tK_e s\right]Q(s)=K_tU(s).
\]

Por tanto:
\begin{equation}
G(s)=\frac{Q(s)}{U(s)}
=
\frac{K_t}{(Ls+R)(ms^2+cs+k)+K_tK_e s}.
\label{eq:Gexact}
\end{equation}

Expandiendo el denominador:
\begin{equation}
G(s)=
\frac{K_t}
{Lm\,s^3 + (Lc+Rm)s^2 + (Lk+Rc+K_tK_e)s + Rk}.
\label{eq:Gexpanded}
\end{equation}

Aqu\'i se ve de forma totalmente expl\'icita de d\'onde salen los cuatro coeficientes del polinomio caracter\'istico.

\section{An\'alisis de estabilidad del modelo lineal exacto}

La ecuaci\'on caracter\'istica es:
\[
Lm\,s^3 + (Lc+Rm)s^2 + (Lk+Rc+K_tK_e)s + Rk = 0.
\]

Si definimos:
\[
a_3=Lm,\qquad
a_2=Lc+Rm,\qquad
a_1=Lk+Rc+K_tK_e,\qquad
a_0=Rk,
\]
entonces, para un polinomio c\'ubico,
\[
a_3 s^3 + a_2 s^2 + a_1 s + a_0,
\]
Routh-Hurwitz exige:
\begin{equation}
a_3>0,\quad a_2>0,\quad a_1>0,\quad a_0>0,\quad a_2a_1>a_3a_0.
\label{eq:RH}
\end{equation}

Como los par\'ametros f\'isicos son positivos, las primeras cuatro condiciones se satisfacen. La condici\'on no trivial es:
\begin{equation}
(Lc+Rm)(Lk+Rc+K_tK_e) > LmRk.
\label{eq:RHcond}
\end{equation}

Esa es la condici\'on de estabilidad del modelo lineal exacto.

\section{C\'omo se obtiene el modelo simplificado que ven\'iamos usando}

El modelo del c\'odigo previo usa un actuador de primer orden:
\begin{equation}
\tau_a \dot F_a + F_a = K_a u.
\label{eq:force1st}
\end{equation}

Ahora explicamos de d\'onde sale.

\subsection{Primera simplificaci\'on: despreciar el back-EMF en la banda de dise\~no}

Si en la banda de inter\'es:
\[
|K_e\dot q| \ll |u|,
\]
entonces \eqref{eq:elec_exact} se aproxima como:
\begin{equation}
L\dot i + Ri \approx u.
\label{eq:nobemf}
\end{equation}

\subsection{Segunda simplificaci\'on: cambiar variable de estado de corriente a fuerza}

Como $F_a=K_t i$, entonces:
\[
\dot F_a = K_t \dot i.
\]

Multiplicando \eqref{eq:nobemf} por $K_t/R$:
\[
\frac{L}{R}\dot F_a + F_a = \frac{K_t}{R}u.
\]

Si definimos:
\begin{equation}
\tau_a = \frac{L}{R},
\qquad
K_a = \frac{K_t}{R},
\label{eq:map}
\end{equation}
obtenemos:
\begin{equation}
\tau_a \dot F_a + F_a = K_a u.
\end{equation}

Esta es exactamente la estructura simplificada usada en el modelo anterior. Entonces:

\begin{quote}
\textbf{el actuador de primer orden no estaba inventado; era una versi\'on simplificada del actuador electromec\'anico cuando se desprecia el back-EMF.}
\end{quote}

\section{Qu\'e se hace primero en la actividad}

La secuencia correcta, con rigor mec\'anico y matem\'atico, es:

\begin{enumerate}[leftmargin=1.2cm]
\item definir el sistema f\'isico y sus variables;
\item escribir las ecuaciones a partir de Newton, Hooke, amortiguamiento viscoso y KVL;
\item construir el modelo realista no lineal;
\item encontrar el equilibrio;
\item linealizar por Taylor/Jacobianos;
\item obtener la funci\'on de transferencia del modelo nominal;
\item analizar estabilidad y dise\~nar el controlador con ese modelo nominal;
\item comparar la respuesta del modelo nominal con la del modelo realista en Simulink/Simscape.
\end{enumerate}

Entonces, respondiendo a tu duda de forma muy precisa:

\begin{quote}
\textbf{no es que Simscape te entregue la ecuaci\'on de tercer orden. Primero se deriva matem\'aticamente la din\'amica; luego Simscape sirve para implementar y contrastar una versi\'on m\'as realista de esa misma planta.}
\end{quote}

\section{Glosario de variables, par\'ametros y s\'imbolos}

En esta secci\'on dejo un glosario de todas las magnitudes importantes que aparecen en las ecuaciones del modelo. La idea es que puedas leer cualquier ecuaci\'on del documento y saber exactamente qu\'e representa cada s\'imbolo.

\subsection{Convenciones de notaci\'on}

\begin{itemize}[leftmargin=1.2cm]
\item Si $r(t)$ es una variable cualquiera, entonces $\dot r = \dfrac{dr}{dt}$ es su primera derivada temporal y $\ddot r = \dfrac{d^2r}{dt^2}$ es su segunda derivada temporal.
\item El s\'imbolo $\tilde{(\cdot)}$ indica \textbf{variable incremental} respecto a un equilibrio est\'atico.
\item El s\'imbolo $\delta(\cdot)$ indica \textbf{peque\~na perturbaci\'on} o desviaci\'on local alrededor del punto de linealizaci\'on.
\item El super\'indice ${}^\star$ indica \textbf{valor de equilibrio} o punto de operaci\'on.
\item Las letras may\'usculas en Laplace, como $Q(s)$ o $U(s)$, representan la transformada de Laplace de sus respectivas variables temporales.
\end{itemize}

\subsection{Variables cinem\'aticas y de salida}

\renewcommand{\arraystretch}{1.2}
\begin{center}
\small
\begin{tabular}{>{\raggedright\arraybackslash}p{2.8cm} >{\raggedright\arraybackslash}p{9.2cm} >{\raggedright\arraybackslash}p{2.2cm}}
\toprule
\textbf{S\'imbolo} & \textbf{Significado} & \textbf{Unidad} \\
\midrule
$t$ & Tiempo. Variable independiente del problema din\'amico. & s \\
$x(t)$ & Posici\'on absoluta de la plataforma superior. & m \\
$z(t)$ & Posici\'on absoluta de la base vibratoria. & m \\
$q(t)$ & Desplazamiento relativo entre plataforma y base, definido como $q=x-z$. Tambi\'en representa la deformaci\'on de la suspensi\'on. & m \\
$\dot x(t)$ & Velocidad absoluta de la plataforma. & m/s \\
$\dot z(t)$ & Velocidad absoluta de la base. & m/s \\
$\dot q(t)$ & Velocidad relativa entre plataforma y base. & m/s \\
$\ddot x(t)$ & Aceleraci\'on absoluta de la plataforma. & m/s$^2$ \\
$\ddot z(t)$ & Aceleraci\'on absoluta de la base. Es la excitaci\'on externa de soporte. & m/s$^2$ \\
$\ddot q(t)$ & Aceleraci\'on relativa entre plataforma y base. & m/s$^2$ \\
$v$ & Variable auxiliar definida como $v=\dot q$. Se usa para expresar el sistema en espacio de estados. & m/s \\
$y$ & Salida del sistema. Seg\'un el caso, puede ser la posici\'on relativa $q$ o la posici\'on absoluta $x=q+z$. & m \\
$w$ & Perturbaci\'on reescrita como aceleraci\'on de base: $w=\ddot z$. & m/s$^2$ \\
$a$ & Aceleraci\'on gen\'erica usada en la forma general de Newton, $\sum F = ma$. & m/s$^2$ \\
\bottomrule
\end{tabular}
\end{center}

\subsection{Fuerzas, se\~nales el\'ectricas y funciones no lineales}

\begin{center}
\small
\begin{tabular}{>{\raggedright\arraybackslash}p{3.0cm} >{\raggedright\arraybackslash}p{9.0cm} >{\raggedright\arraybackslash}p{2.2cm}}
\toprule
\textbf{S\'imbolo} & \textbf{Significado} & \textbf{Unidad} \\
\midrule
$\sum F$ & Suma de todas las fuerzas aplicadas sobre la masa de la plataforma. & N \\
$F_s$ & Fuerza del resorte lineal, dada por la ley de Hooke. & N \\
$F_d$ & Fuerza del amortiguador viscoso. & N \\
$F_a$ & Fuerza activa aplicada por el actuador entre base y plataforma. & N \\
$F_{\text{pert}}$ & Fuerza inercial equivalente asociada a la excitaci\'on de base; en coordenada relativa vale $-m\ddot z$. & N \\
$F_{s,\text{nl}}$ & Fuerza el\'astica no lineal cuando se incluye el t\'ermino c\'ubico. & N \\
$F_f$ & Fuerza de fricci\'on seca regularizada. & N \\
$F_{\text{topes}}$ & Fuerza de contacto generada cuando la deformaci\'on excede el recorrido permitido. & N \\
$\Delta \ell$ & Deformaci\'on geom\'etrica del resorte respecto a su longitud de equilibrio. En este sistema, $\Delta \ell=x-z=q$. & m \\
$\Delta \dot \ell$ & Velocidad relativa de deformaci\'on del amortiguador. En este sistema, $\Delta \dot \ell=\dot x-\dot z=\dot q$. & m/s \\
$u(t)$ & Voltaje de mando o se\~nal de control aplicada al actuador. & V \\
$u_{\text{act}}$ & Voltaje efectivamente aplicado luego de pasar por la saturaci\'on del actuador. & V \\
$\operatorname{sat}(u)$ & Funci\'on de saturaci\'on. Limita la se\~nal de control a un rango f\'isicamente realizable. & V \\
$i(t)$ & Corriente el\'ectrica del actuador. Es el tercer estado del modelo electrome\'canico exacto. & A \\
$e_b$ & Fuerza contraelectromotriz o back-EMF inducida por el movimiento relativo del actuador. & V \\
$\tanh(\cdot)$ & Funci\'on tangente hiperb\'olica usada para suavizar la fricci\'on seca y hacerla derivable en velocidad cero. & adim. \\
$\operatorname{sech}(\cdot)$ & Funci\'on hiperb\'olica que aparece al derivar $\tanh(\cdot)$ durante la linealizaci\'on. & adim. \\
\bottomrule
\end{tabular}
\end{center}

\subsection{Par\'ametros f\'isicos del sistema}

\begin{center}
\small
\begin{tabular}{>{\raggedright\arraybackslash}p{3.0cm} >{\raggedright\arraybackslash}p{9.0cm} >{\raggedright\arraybackslash}p{2.2cm}}
\toprule
\textbf{S\'imbolo} & \textbf{Significado} & \textbf{Unidad} \\
\midrule
$m$ & Masa equivalente de la plataforma superior y del equipo sensible montado sobre ella. & kg \\
$k$ & Rigidez lineal equivalente del elemento el\'astico entre base y plataforma. & N/m \\
$c$ & Coeficiente de amortiguamiento viscoso equivalente entre base y plataforma. & N\,s/m \\
$k_3$ & Coeficiente de rigidez no lineal c\'ubica. Modela desviaciones de la ley de Hooke para deformaciones mayores. & N/m$^3$ \\
$F_c$ & Nivel o amplitud de la fricci\'on tipo Coulomb en la regularizaci\'on suave. & N \\
$v_\varepsilon$ & Velocidad de suavizado de la fricci\'on regularizada. Controla qu\'e tan abrupta es la transici\'on de $\tanh(\dot q/v_\varepsilon)$. & m/s \\
$L$ & Inductancia del devanado del actuador. Introduce la din\'amica el\'ectrica de la corriente. & H \\
$R$ & Resistencia el\'ectrica del devanado del actuador. & $\Omega$ \\
$K_t$ & Constante de fuerza del actuador. Relaciona fuerza con corriente: $F_a=K_t i$. & N/A \\
$K_e$ & Constante de back-EMF. Relaciona voltaje inducido con velocidad relativa: $e_b=K_e\dot q$. & V\,s/m \\
$q_{\max}$ & M\'axima deformaci\'on admisible antes de que entren en contacto los topes mec\'anicos. & m \\
$k_t$ & Rigidez equivalente del tope mec\'anico cuando hay contacto. & N/m \\
$c_t$ & Amortiguamiento equivalente del contacto con el tope. & N\,s/m \\
$g$ & Aceleraci\'on de la gravedad. Solo aparece expl\'icitamente si se modela el movimiento vertical antes de trasladarse al equilibrio est\'atico. & m/s$^2$ \\
$mg$ & Peso de la plataforma. Es la fuerza gravitacional aplicada sobre la masa. & N \\
$\tau_a$ & Constante de tiempo del actuador simplificado de primer orden, definida como $\tau_a=L/R$. & s \\
$K_a$ & Ganancia est\'atica del actuador simplificado en la relaci\'on entre voltaje y fuerza, definida como $K_a=K_t/R$. & N/V \\
\bottomrule
\end{tabular}
\end{center}

\subsection{Variables de equilibrio, linealizaci\'on y modelo incremental}

\begin{center}
\small
\begin{tabular}{>{\raggedright\arraybackslash}p{3.0cm} >{\raggedright\arraybackslash}p{9.0cm} >{\raggedright\arraybackslash}p{2.2cm}}
\toprule
\textbf{S\'imbolo} & \textbf{Significado} & \textbf{Unidad} \\
\midrule
$x_0$ & Posici\'on absoluta de equilibrio est\'atico de la plataforma cuando se considera gravedad en eje vertical. & m \\
$z_0$ & Posici\'on absoluta de equilibrio est\'atico de la base. & m \\
$F_{a0}$ & Fuerza del actuador en el equilibrio est\'atico. & N \\
$\tilde x$ & Desviaci\'on incremental de la posici\'on de la plataforma respecto al equilibrio est\'atico. & m \\
$\tilde z$ & Desviaci\'on incremental de la posici\'on de la base respecto al equilibrio est\'atico. & m \\
$\tilde F_a$ & Desviaci\'on incremental de la fuerza del actuador respecto al equilibrio est\'atico. & N \\
$x^\star$ & Vector de estado en el punto de equilibrio usado para linealizar. & seg\'un componente \\
$u^\star$ & Entrada de equilibrio usada para linealizar. & V \\
$w^\star$ & Perturbaci\'on de equilibrio usada para linealizar. & m/s$^2$ \\
$\delta x$ & Peque\~na variaci\'on del vector de estado alrededor del equilibrio: $\delta x=x-x^\star$. & seg\'un componente \\
$\delta u$ & Peque\~na variaci\'on de la entrada alrededor del equilibrio. & V \\
$\delta w$ & Peque\~na variaci\'on de la perturbaci\'on alrededor del equilibrio. & m/s$^2$ \\
$c_{\text{eq}}$ & Amortiguamiento equivalente local que aparece al linealizar estrictamente la fricci\'on regularizada: $c_{\text{eq}}=c+F_c/v_\varepsilon$. & N\,s/m \\
$\xi$ & Variable auxiliar adimensional usada en la aproximaci\'on local $\tanh(\xi)\approx \xi$. & adim. \\
\bottomrule
\end{tabular}
\end{center}

\subsection{Estados, matrices y dominio de Laplace}

\begin{center}
\small
\begin{tabular}{>{\raggedright\arraybackslash}p{3.0cm} >{\raggedright\arraybackslash}p{9.0cm} >{\raggedright\arraybackslash}p{2.2cm}}
\toprule
\textbf{S\'imbolo} & \textbf{Significado} & \textbf{Unidad} \\
\midrule
$x_1$ & Primer estado del sistema, igual a la deformaci\'on relativa: $x_1=q$. & m \\
$x_2$ & Segundo estado del sistema, igual a la velocidad relativa: $x_2=\dot q=v$. & m/s \\
$x_3$ & Tercer estado del sistema, igual a la corriente del actuador: $x_3=i$. & A \\
$x$ & Vector de estado completo: $x=[x_1\;x_2\;x_3]^T$. & mixta \\
$A$ & Matriz de din\'amica del sistema linealizado. & depende del elemento \\
$B$ & Matriz de entrada del sistema linealizado respecto a la se\~nal de control. & depende del elemento \\
$E$ & Matriz de entrada de perturbaci\'on del sistema linealizado respecto a $w=\ddot z$. & depende del elemento \\
$C$ & Matriz de salida del sistema cuando se selecciona $y=q$ o la variable de salida correspondiente. & depende del elemento \\
$s$ & Variable compleja del dominio de Laplace. & 1/s \\
$Q(s)$ & Transformada de Laplace de $q(t)$. & m \\
$I(s)$ & Transformada de Laplace de $i(t)$. & A \\
$U(s)$ & Transformada de Laplace de $u(t)$. & V \\
$G(s)$ & Funci\'on de transferencia de la planta nominal. En este proyecto, relaciona $Q(s)$ con $U(s)$. & m/V \\
$a_3,a_2,a_1,a_0$ & Coeficientes del polinomio caracter\'istico c\'ubico de la planta lineal exacta. Se usan en el criterio de Routh-Hurwitz. & seg\'un coeficiente \\
\bottomrule
\end{tabular}
\end{center}

\subsection{Lectura r\'apida del glosario}

Si tienes que defender el modelo oralmente, lo m\'as importante es no confundir estas cuatro cosas:

\begin{itemize}[leftmargin=1.2cm]
\item $x$ es posici\'on absoluta de la plataforma, mientras que $q$ es desplazamiento relativo.
\item $u$ es voltaje de control, mientras que $i$ es corriente y $F_a$ es fuerza.
\item $w=\ddot z$ no es una fuerza aplicada directamente, sino la aceleraci\'on de la base reescrita como perturbaci\'on.
\item $A$, $B$, $C$ y $E$ no son variables f\'isicas medibles por s\'i solas, sino matrices del modelo linealizado.
\end{itemize}

\section{Conclusi\'on}

La formulaci\'on rigurosa y defendible del proyecto queda as\'i:

\begin{itemize}[leftmargin=1.2cm]
\item la ecuaci\'on mec\'anica sale de la segunda ley de Newton;
\item el resorte sale de la ley de Hooke;
\item el amortiguador sale del modelo viscoso ideal;
\item el actuador sale de fuerza proporcional a corriente m\'as KVL y back-EMF;
\item el sistema nominal exacto es de tercer orden;
\item el modelo realista agrega no linealidades y restricciones f\'isicas;
\item la linealizaci\'on se obtiene por Taylor de primer orden alrededor del equilibrio;
\item las expresiones simplificadas aparecen en pasos identificables y justificables.
\end{itemize}

Si te preguntan en la exposici\'on ``por qu\'e eliminaste tal t\'ermino'', la respuesta correcta ya no ser\'a ``porque se desprecia'', sino:

\begin{quote}
\textit{porque al evaluar el Jacobiano en el equilibrio ese t\'ermino no aporta en primer orden}, o bien \textit{porque se reserv\'o para el modelo realista y no para el nominal de dise\~no}.
\end{quote}

\end{document}
